\subsection{Scalar relaxation}
\label{sec:relaxation-convergence}

Relaxed-SAV schemes modify the auxiliary state while preserving a discrete
energy bound \cite{jiang2022rsav,zhang2022gsavrelax,huang2024optimalrsav}.
The squared-scalar update below is the relaxation of \cite{zhang2026pbsav};
it depends on the update only through its Lyapunov functional, so one
statement covers both updates of this section.

Without relaxation, \(q^{n+1}=q^n+d^{n+1}\) is the stored scalar in
\cref{eq:direct-family,eq:abstract-momentum-states}. When relaxation is
applied, the same algebraic output is provisional and is denoted by
\begin{equation}
  \bar q^{n+1}=q^n+d^{n+1}.
  \label{eq:relaxation-candidate}
\end{equation}
For the direct update the old and provisional energies are
\(\Hcal_{\mathrm{dir}}^n=(q^n)^2\) and
\(\bar\Hcal_{\mathrm{dir}}^{n+1}=(\bar q^{n+1})^2\); for the momentum update
they are \(\Hcal_{\mathrm{mom}}^n=(q^n)^2+\norm{p^n}_{M^n}^2/(2\eta)\) and
\(\bar\Hcal_{\mathrm{mom}}^{n+1}=(\bar q^{n+1})^2
+\norm{p^{n+1}}_{M^{n+1}}^2/(2\eta)\).

\begin{proposition}
\label{prop:generic-relaxation}
For either update, let \((\Hcal^n,\bar\Hcal^{n+1})\) denote the corresponding
pair above and suppose that
\begin{equation}
  \Hcal^n-\bar\Hcal^{n+1}=\mathcal D^n\geq0.
  \label{eq:generic-unrelaxed-dissipation}
\end{equation}
For \(0\leq\rho^n\leq1\), define
\begin{equation}
  (q^{n+1})^2
  =\min\left\{
  Q(\theta^{n+1})^2,\;
  (\bar q^{n+1})^2+\rho^n\mathcal D^n
  \right\},
  \qquad q^{n+1}\geq0.
  \label{eq:generic-relaxation}
\end{equation}
Let \(\Hcal^{n+1}\) be the same energy evaluated with the accepted
scalar \(q^{n+1}\). Then
\begin{equation}
  \Hcal^{n+1}
  \leq\Hcal^n-(1-\rho^n)\mathcal D^n
  \leq\Hcal^n.
  \label{eq:relaxed-energy-bound}
\end{equation}
\end{proposition}

\begin{proof}
Relaxation changes only the scalar part of the new energy.  Therefore
\(
\Hcal^{n+1}\leq\bar\Hcal^{n+1}+\rho^n\mathcal D^n
\), and \cref{eq:generic-unrelaxed-dissipation} gives
\cref{eq:relaxed-energy-bound}.
\end{proof}

With \(q^0=Q(\theta^0)\), the accepted scalars satisfy
\(0\leq q^n\leq Q(\theta^n)\) by \cref{eq:generic-relaxation}.
Full relaxation \(\rho^n=1\) preserves the modified energy bound; choosing
\(\rho^n\leq\bar\rho<1\) also gives
\(\sum_n\mathcal D^n\leq\Hcal^0/(1-\bar\rho)\).
Under the curvature condition in \ref{app:directional-coverage}, full
relaxation preserves \(q^n=Q(\theta^n)\), so the momentum energy becomes
\(F(\theta^n)+C+\norm{p^n}_{M^n}^2/(2\eta)\).
